\mwnewpart
\renewcommand\runningtitle{Section 5. The IBM Commands}
\chapter{The IBM Commands}
\sectioncontents

\section{Writing Them Again}\ixs{IBM commands}

\begin{plainwords}
IBM wrote a dozen of PC DOS 2.10's commands itself, and never published
them. RetrOS/LT has its own, written to behave as IBM's do: the same
switches, the same questions, the same messages, and where it matters to
other programs, the same layout in memory and on disk. They are small
programs, most under 5,000 bytes, written for NASM, and they use only
what an 8088 and the PC's BIOS provide.
\end{plainwords}

Each was written from IBM's documentation of PC DOS 2.10 and what is
known of its behavior, and tested against the situations the
documentation describes. Where IBM's behavior could not be learned, the
choice made is noted.

\section{FORMAT: the IBM Part}\ixs{FORMAT command}

Microsoft's \typed{FORMAT} leaves everything that depends on the computer
to an OEM module, through a fixed interface (\file{FORMAT.txt} of the
kit): \routine{INIT}, \routine{DISKFORMAT}, \routine{BADSECTOR},
\routine{WRTFAT} and \routine{DONE}, with \code{SWITCHLIST},
\code{FATID}, \code{STARTSECTOR} and the rest. \file{src/IBMFOR.ASM} is
RetrOS/LT's, linked after \file{FORMAT.ASM} and \file{FORMES.ASM}.

\begin{itemize}
\item \textbf{Diskettes} are formatted a track at a time through the ROM
  (interrupt 13h, function 5) and verified; a bad track is remembered and
  reported to \routine{BADSECTOR}. \code{SWITCHLIST} is IBM's,
  \code{"B81OVS"}.
\item \textbf{A fixed disk} is not low-level formatted. Its layout comes
  from the drive parameter block \file{RETBIO.COM} built (function 32h),
  its geometry from the ROM, and where its partition starts from the
  master boot record; every sector is then verified, a track at a time,
  and bad ones reported.
\item \textbf{WRTFAT} writes the startup record and the file allocation
  tables straight to the disk: until the FAT ID byte is there, the kernel
  cannot know what the disk is.
\item \textbf{The messages} \shown{Formatting...} and \shown{Format
  complete}, which PC DOS 2's \typed{FORMAT} shows and Microsoft's leaves to
  the OEM part, come from \routine{DISKFORMAT} and the first
  \routine{WRTFAT}.
\end{itemize}

\typed{FORMAT} calls the ROM itself, so it must keep its own transfers off
64 KB boundaries, as \file{RETBIO.COM} does:

\srclist{ibmfor_bounce}{FORMAT's bounce buffer.}

\engnote{The first version did not, and \typed{FORMAT} failed with
\shown{Format failure} whenever DOS happened to load it where its startup
record buffer crossed a 64 KB boundary. The review found it; a resident
program of the right size reproduced it (Section~7). \typed{DISKCOPY},
\typed{DISKCOMP}, \typed{FDISK} and Install had the same fault in their
own direct calls, and have the same remedy.}

\section{SYS}\ixs{SYS command}

Microsoft's \typed{SYS} copies the system files and, for the IBM version,
writes a startup record. \patch{0005} makes it always write RetrOS/LT's:
it reads the disk's BIOS parameter block (interrupt 25h) and keeps it, or,
on a diskette with none, builds a DOS 1 one from the FAT ID byte, and
writes the RetrOS/LT code around it.

\section{FDISK}\ixs{FDISK command}

\typed{FDISK} has PC DOS 2.10's menu: create the DOS partition, change the
active partition, delete the DOS partition, display the partition table,
and, with two fixed disks, select the other. Creating offers the whole
disk, which gives the largest free run of cylinders, to at most 65,535
sectors; or a size and a first cylinder, checked against the other
partitions. A partition on cylinder 0 starts on head 1, after the master
boot record's track. A new disk with no master boot record gets RetrOS/LT's
(\file{MBRCODE.INC}). On leaving after a change, \typed{FDISK} restarts the
computer through interrupt 19h, so that \file{RETBIO.COM} finds the new
partition.

\section{DISKCOPY and DISKCOMP}\ixs{DISKCOPY command}\ixs{DISKCOMP command}

Both work through the ROM, a track at a time, with as many tracks in
memory at once as fit, each placed so that it does not cross a 64 KB
boundary; they share \file{cmd/FLOPPY.INC}. The source's format comes from
its FAT ID byte, and the diskette parameter table is given its sectors per
track. \typed{DISKCOPY} formats a target track that will not take the data
as it is, and says \shown{Formatting while copying} the first time.
\typed{DISKCOMP} reports a difference by side and track. Both take
\typed{/1}, \typed{DISKCOMP} \typed{/8} too.

\section{COMP and TREE}\ixs{COMP command}\ixs{TREE command}

\typed{COMP} compares files byte for byte, ten differences at most, with
IBM's prompts when a name is left out, wildcards in the first name and a
template in the second, and \shown{Compare more files (Y/N)?} at the end.
\typed{TREE} walks the directories depth first, with a disk transfer area
for each level, and prints IBM's \shown{Path:}, \shown{Sub-directories:}
and \shown{Files:} lines.

\section{ASSIGN}\ixs{ASSIGN command}

\typed{ASSIGN} stays resident, catching interrupt 21h (functions that
name a drive in \code{DL}, in a file control block, or in a path),
interrupts 25h and 26h. PC DOS 2's \typed{FORMAT} looks for it, to refuse
to work on an assigned drive, by a fixed layout: the drive table at offset
103h of its segment and its interrupt 21h handler at 122h. RetrOS/LT's
\typed{ASSIGN} is laid out the same way:

\srclist{assign_layout}{ASSIGN's layout.}

A path's drive letter is changed where it is, in the caller's string, and
the request passed on:

\srclist{assign_path}{A path request.}

\engnote{Two other ways were tried. Calling DOS from the handler and
returning broke EXEC, which in DOS 2 does not come back to its caller in
the ordinary way. Changing a copy of the name, and giving DOS the copy, kept
the caller's string unchanged, but \file{COMMAND.COM} works out
directories from the string it passed: with the copy, \typed{REN} on an
assigned drive renamed the first file of the directory,
\file{COMMAND.COM}, instead of the one named. Changing the string where it
is is what \file{COMMAND.COM} expects.}

\section{MODE}\ixs{MODE command}

\typed{MODE} sets printer width and line spacing by sending the IBM
printer's control codes, the serial port with interrupt 14h, and the
display with interrupt 10h and the equipment flags at 0040:0010. For
rerouting (\typed{LPT1:=COM1}), retry (\typed{P}) and shifting the
display, a resident part hooks interrupts 17h and 14h, and copies the
ROM's video parameter table (interrupt 1Dh) so it can change the
horizontal sync position. A later \typed{MODE} finds the resident part by
the word \code{MODE} just before its handler:

\srclist{mode_int17}{MODE's printer handler: rerouting and retry.}

Every parameter is read and checked before anything is done, so a mistake
changes nothing.

\section{GRAPHICS}\ixs{GRAPHICS command}

\typed{GRAPHICS} takes interrupt 5, the print-screen interrupt. In a text
mode it passes the request to the ROM. In modes 4, 5 and 6 it prints the
screen from video memory at B800:0000 on the IBM Graphics Printer, turned
on its side: each printer line is a column of eight screen dots, printed
from the bottom row up with \code{ESC K} graphics, line spacing 8/72
inch. A 640 by 200 screen prints each dot twice across; a 320 by 200
screen prints each dot as two by two printer dots, the four colors as
none, one, two and four of them. The status byte at 0050:0000 is kept as
the ROM keeps it.

\section{BACKUP and RESTORE}\ixs{BACKUP command}\ixs{RESTORE command}

The diskettes are in PC DOS 2's format, so that the two systems can read
each other's backups:

\srclist{br_format}{The backup format.}

\engnote{Whether DOS 2's \typed{BACKUP} kept the files in the root
directory or in directories mirroring the source could not be settled from
the documents found. RetrOS/LT keeps them in the root, which the header's
full path and the 3.x documents' mention of a pack switch that makes
directories only when the root is full both point to. A file whose name is
already on the diskette, from another directory, goes on the next
diskette.}

\typed{BACKUP} erases the target's root directory files (unless
\typed{/A}), writes \file{BACKUPID.@@@}, copies each file behind its
header with its date and time, splits a file across diskettes when the
space runs out, clears the archive bit, and marks the last diskette.
\typed{RESTORE} reads the diskettes in order, matches each file's path
against the one asked for, makes directories, follows split files, and
with \typed{/P} asks before replacing a read-only or changed file. Both
set the exit codes 0, 1, 3 and 4.

\section{ANSI.SYS}\ixs{ANSI.SYS}

\file{ANSI.SYS} is a character device named \code{CON}, which therefore
takes \code{CON}'s place. It writes through the ROM with an attribute of
its own, keeps a parameter list for each escape sequence, and has a queue
for what it gives programs to read: a reassigned key's text, the second
byte of an extended key, and the cursor report. It takes interrupt 29h and
interrupt 1Bh itself. A key is looked up in the reassignment table as it
is read:

\srclist{ansi_keyin}{A key, and its reassignment.}

\engnote{The first version passed \keys{Ctrl}{PrtSc}, which the ROM gives
as 7200h, through as an extended key, so that with \file{ANSI.SYS} loaded
it no longer turned printer echo on and off. The review found it; it is
now \keys{Ctrl}{P}, as \file{RETBIO.COM}'s \code{CON} gives it.}
